Nelle ultime decadi, la valutazione del comportamento sismico degli edifici è notevolmente progredita rispetto ai tradizionali obiettivi riguardanti principalmente la salvaguardia della vita. L’impatto economico e l’interruzione di attività nevralgiche a seguito dei considerevoli danni strutturali causati dal terremoto di Northridge negli Stati Uniti nel 1994 hanno evidenziato la necessità di un cambiamento nelle modalità di definizione delle prestazioni degli edifici. Le prestazioni esibite dagli edifici durante il sisma che ha colpito L’Aquila nel 2009, città ancora in fase di ricostruzione, e i report iniziali dei danni a seguito del sisma di Rieti nel 2016 hanno confermato l’esigenza di una valutazione del comportamento più ampia, non esclusivamente legata alla salvaguardia della vita degli occupanti. Questa valutazione integrata che considera anche le perdite economiche potenziali dovute al danneggiamento sia degli elementi strutturali che degli elementi non-strutturali è conosciuta come metodologia PEER PBEE. In questo modo le “performance” strutturali possono essere meglio descritte da strumenti come le perdite economiche dirette, attirando così l’interesse di un pubblico più vasto rispetto alle sole comunità di ingegneri e progettisti sismici.

Questa tesi ha come oggetto di studio la valutazione del comportamento sismico di edifici a telaio in cemento armato con tamponature di laterizio in Italia progettati esclusivamente per carichi gravatici fino all’introduzione delle prime norme antisismiche nel 1970. L’approccio allo studio di questa tipologia strutturale segue le linee guida suggerite da FEMA P58 il cui obiettivo è quello di fornire agli ingegneri progettisti, attraverso strumenti semplificati e avanzati, la possibilità di caratterizzare in maniera significativa il comportamento degli edifici e interagire in maniera più efficace con coloro che prendono decisioni. Il lavoro ha previsto la calibrazione di modelli numerici in OpenSees di questi telai, considerando non solo la non-linearità degli elementi critici della struttura ma anche la perdita progressiva di capacità e rigidezza caratterizzanti il comportamento a collasso di questa tipologia strutturale. Ulteriori strumenti necessari ad un’adeguata valutazione delle perdite sono stati indagati e implementati in una MATLAB routine, come le implicazioni legate alla risposta torsionale delle strutture per esempio. 

Un dettagliato studio su questi modelli numerici, calibrati attraverso le prove sperimentali disponibili in letteratura, è stato condotto per poter quantificare l’incertezza legata ai vari parametri di modellazione. É stato possibile in questo modo considerare l’effetto di tale incertezza sulla dispersione globale della domanda e l’impatto sulle funzioni di fragilità a collasso di questi edifici. Valori empirici di incertezza legata alla modellazione sono stati proposti per telai di calcestruzzo armato con tamponature. Comparando i risultati ottenuti è stato possibile osservare come i parametri proposti da FEMA P58 non possono essere utilizzati per la tipologia strutturale oggetto di studio. Le analisi realizzate su un numero di modelli tale da caratterizzare pienamente la risposta di questa tipologia strutturale in termini sia di domanda attesa che di capacità ultima, hanno evidenziato come considerare il potenziale collasso a taglio sui pilastri, dovuto all’interazione con la tamponatura, porta a una riduzione del 20\% dell’intensità mediana a collasso. Il lavoro ha previsto inoltre non solo il calcolo delle perdite attese per le tipologie di edifici analizzati ma anche l’impatto positivo su perdite attese e performance strutturali legato ad efficaci interventi di miglioramento sismico. Si è evidenziato infatti, quando la sicurezza, in zone di bassa o moderata sismicità in Italia, le risposta in termini di EAL può essere notevolmente ridotta attraverso abili strategie di adeguamento sismico dei più vulnerabili elementi non strutturali. Al contrario si è evidenziato come strategie basate sull’irrigidimento e l’incremento della capacità ultima secondo le disposizioni NTC 2008 possano condurre ad un peggioramento della risposta strutturale in termini di EAL. Ciò è legato alla relazione fra ridotto spostamento interpiano e incremento dell’accelerazione di piano stessa delle strutture adeguate oltre al potenziale incremento dell’intensità di misura considerata dato un particolare periodo di ritorno e secondo l’intensità di misura considerata. Questo è stato confermato dalle osservazioni post-sisma seguenti l’evento verificatosi in Ecuador nel 2016, dove ospedali che avevano subito interventi di adeguamento basati sull’incremento di rigidezza e capacità ultima hanno evidenziato estesi danni a componenti non strutturali e ad apparecchiature cruciali inficiando la funzionalità e l’agibilità delle strutture stesse immediatamente dopo il sisma.

Il lavoro si inserisce nell’ambito delle metodologie di valutazione sismica proposte da FEMA P58, dove metodologie di analisi dettagliate descritte precedentemente sono affiancate a metodologie semplificate che non richiedono analisi numeriche come il DBA. La metodologia DBA è stata confrontata con metodologie di analisi avanzate per mostrare come, nonostante il metodo sembri funzionare abbastanza bene, ulteriori elementi debbano essere investigati prima di formulare una proposta finale di DBA per questa tipologia di edifici. In ogni caso lo studio svolto consente di calcolare la domanda mediana sull’edificio attraverso la metodologia DBA e la dispersione legata a tale domanda oltre che la correlazione fra i vari EDP attraverso ulteriori approfondimenti contenuti nella tesi. Il confronto tra metodologie di approccio avanzate e DBA applicate ad un edificio reale adibito a scuola ha mostrato come i valori empirici di dispersione proposti e la correlazione fra i vari parametri di domanda possono rappresentare in maniera soddisfacente la risposta di questi edifici. I valori proposti da FEMA P58 non possono essere considerati soddisfacenti in quanto sottostimano la dispersione della domanda, sovrastimano la correlazione fra i parametri di domanda sottostimando conseguentemente le perdite attese e la loro variabilità. 
Infine la metodologia SAC/FEMA per il calcolo della la probabilità media annua di superamento di differenti stati limite, è stata sviluppata per questa tipologia strutturale, mostrando come il termine in potenza b nella relazione domanda-intensità, tipicamente assunto uguale a 1.0, non è rappresentativo di questa tipologia di telai. Valori più appropriati di b sono stati identificati per telai con tamponature.
